\documentclass[11pt,a4paper]{article}
\usepackage[italian]{babel}
\usepackage[margin=2.5cm]{geometry}
\usepackage{parskip}
\usepackage{amsmath, amssymb, amsthm}
\usepackage{booktabs}
\usepackage{array}
\usepackage{xcolor}
\usepackage{needspace}
\usepackage{enumitem}
\usepackage{listings}
\usepackage{tikz}
\usepackage[most]{tcolorbox}
\usepackage[hidelinks]{hyperref}
\usetikzlibrary{decorations.pathreplacing, shapes.geometric, arrows.meta, positioning}

% Niente vedove/orfane: i paragrafi non lasciano righe isolate a fine/inizio pagina
\widowpenalty=10000
\clubpenalty=10000

% Elenchi più compatti
\setlist{itemsep=2pt, parsep=0pt, topsep=4pt, partopsep=0pt}

% --- Riquadri con bordo nero, senza numero (si spezzano tra due pagine) ---
% Uso: \begin{definizione} ... \end{definizione}
%      \begin{teorema}[Nome] ... \end{teorema}  (il nome è facoltativo)
\tcbset{nota/.style={enhanced jigsaw, breakable, colback=white,
  colframe=black, boxrule=0.5pt, sharp corners}}
\NewTColorBox{definizione}{}{nota,
  before upper={\textbf{Definizione.}\ }}
\NewTColorBox{teorema}{o}{nota,
  before upper={\textbf{Teorema\IfValueT{#1}{ (#1)}.}\ }}
\NewTColorBox{proposizione}{o}{nota,
  before upper={\textbf{Proposizione\IfValueT{#1}{ (#1)}.}\ }}
\NewTColorBox{proprieta}{o}{nota,
  before upper={\textbf{Proprietà\IfValueT{#1}{ (#1)}.}\ }}
\NewTColorBox{assioma}{o}{nota, boxrule=1pt,
  before upper={\textbf{Assioma\IfValueT{#1}{ (#1)}.}\ }}

% --- Ambienti semplici (senza riquadro) ---
% Il titolo sta in cima al blocco, anche se il contenuto inizia con un riquadro o
% con codice; e non resta mai da solo in fondo a una pagina.
% Uso: \begin{esempio}[Titolo facoltativo] ... \end{esempio}
\newcommand{\newrunin}[2]{%
  \NewDocumentEnvironment{#1}{o}{%
    \par\Needspace{4\baselineskip}\addvspace{\medskipamount}\noindent
    \textbf{#2}\IfValueT{##1}{ (##1)}\textbf{.}\ \ignorespaces}%
  {\par\addvspace{\medskipamount}}}
\newrunin{esempio}{Esempio}
\newrunin{esercizio}{Esercizio}
\newrunin{osservazione}{Osservazione}

% V e F nelle tavole di verità
\newcommand{\V}{\textsf{V}}
\newcommand{\F}{\textsf{F}}

% --- Codice C ---
% Uso: \begin{codice} ... \end{codice}      codice C numerato
%      \begin{comando} ... \end{comando}    comandi da terminale
%      \begin{risultato} ... \end{risultato} output di un programma
%      \lstinline|printf("ciao");|           codice dentro al testo
\lstdefinestyle{c}{
  language=C,
  basicstyle=\ttfamily\small,
  keywordstyle=\color{blue}\bfseries,
  commentstyle=\color{gray}\itshape,
  stringstyle=\color{red!70!black},
  numbers=left,
  numberstyle=\tiny\color{gray},
  numbersep=8pt,
  columns=flexible,
  keepspaces=true,
  breaklines=true,
  showstringspaces=false,
  tabsize=4,
  morekeywords={include, define},
  % lettere accentate nei commenti
  literate={à}{{\`a}}1 {è}{{\`e}}1 {é}{{\'e}}1 {ì}{{\`i}}1
           {ò}{{\`o}}1 {ù}{{\`u}}1 {È}{{\`E}}1 {À}{{\`A}}1
}
\lstset{style=c}
\newtcblisting{codice}{listing only, nota, left=6mm,
  listing options={style=c}}
\newtcblisting{comando}{listing only, nota,
  listing options={basicstyle=\ttfamily\small, numbers=none, language={},
    columns=flexible, keepspaces=true, breaklines=true}}
\newtcblisting{risultato}{listing only, enhanced jigsaw, breakable,
  colback=black!5, colframe=black, boxrule=0.5pt, sharp corners,
  listing options={basicstyle=\ttfamily\small, numbers=none, language={},
    columns=flexible, keepspaces=true, breaklines=true}}

% --- Diagrammi di flusso (simboli standard) ---
\tikzset{
  terminale/.style = {ellipse, draw, minimum width=2.5cm, minimum height=0.9cm},
  io/.style        = {trapezium, trapezium left angle=70, trapezium right angle=110,
                      draw, minimum width=2.5cm, minimum height=0.9cm},
  processo/.style  = {rectangle, draw, minimum width=2.5cm, minimum height=0.9cm},
  decisione/.style = {diamond, draw, aspect=2, minimum width=2.5cm},
  freccia/.style   = {thick, -{Stealth}}
}

\title{Discontinuità}
\author{Marco Cerbella}
\date{Lezione 10 -- 26 ottobre 2026}

\begin{document}
\maketitle

\section{Discontinuità}

Una funzione non continua in un punto $x_0$ si dice \emph{discontinua} in $x_0$.

\begin{definizione}[Tipi di discontinuità]
\begin{itemize}
    \item \emph{Discontinuità eliminabile}: esiste finito il limite $\displaystyle\lim_{x\to x_0} f(x) \neq f(x_0)$.
    \item \emph{Discontinuità a salto (o di prima specie)}: esistono finiti i limiti $\displaystyle\lim_{x\to x_0^+} f(x) = \ell_1$ e $\displaystyle\lim_{x\to x_0^-} f(x) = \ell_2$, ma sono diversi fra loro. Il \emph{salto} in $x_0$ è $\ell_1 - \ell_2$.
    \item \emph{Discontinuità di seconda specie}: uno dei due limiti $\displaystyle\lim_{x\to x_0^+} f(x)$, $\displaystyle\lim_{x\to x_0^-} f(x)$ non esiste oppure è infinito.
\end{itemize}
\end{definizione}

\section{Esempi}

\begin{esempio}[Discontinuità eliminabile]
Sia
\[
g(x) = \begin{cases} x & x \neq 0 \\ 1 & x = 0 .\end{cases}
\]
Si ha $\displaystyle\lim_{x\to 0} g(x) = \lim_{x\to 0} x = 0 \neq g(0) = 1$, quindi $g$ non è continua in $x_0 = 0$: ha una discontinuità eliminabile.

Cambiando $g$ in un punto, ponendo $\bar g(0) = \lim_{x\to 0} g(x) = 0$, cioè
\[
\bar g(x) = \begin{cases} x & x \neq 0 \\ 0 & x = 0 ,\end{cases}
\]
si ha $\bar g(x) = g(x)$ tranne che in $x = 0$ e $\bar g$ è continua.
\end{esempio}

\begin{esempio}[Funzione mantissa]
La funzione \emph{mantissa} è $M(x) = x - [x]$, dove $[x] = m$ è l'intero tale che $m \le x < m + 1$, con $x \in \mathbb{R}$, $m \in \mathbb{Z}$. Quindi:
\[
\text{se } 0 \le x < 1 \Rightarrow [x] = 0, \qquad \text{se } 1 \le x < 2 \Rightarrow [x] = 1, \dots
\qquad
M(x) = \begin{cases} x & 0 \le x < 1 \\ x - 1 & 1 \le x < 2 \\ x - 2 & 2 \le x < 3 \\ \dots \end{cases}
\]
Il grafico è una successione di segmenti paralleli alla bisettrice, ciascuno da $0$ (incluso) fino a $1$ (escluso).

Consideriamo $M$ con $x \in\, ]0, 2[$:
\[
M(x)\big|_{]0,2[} = \begin{cases} x & 0 < x < 1 \\ x - 1 & 1 \le x < 2 .\end{cases}
\]
Studiamo $x_0 = 1$:
\begin{itemize}
    \item $\displaystyle\lim_{x\to 1^-} M(x) = \lim_{x\to 1^-} x = 1$ e $M(1) = 1 - 1 = 0$, quindi $\lim_{x\to 1^-} M(x) \neq M(1)$: esiste il limite sinistro ma $M$ non è continua da sinistra;
    \item $\displaystyle\lim_{x\to 1^+} M(x) = \lim_{x\to 1^+} (x-1) = 0 = M(1)$: esiste il limite destro ed è uguale a $M(1)$, quindi $M$ è continua da destra;
    \item $\displaystyle\nexists \lim_{x\to 1} M(x)$ perché il limite destro e il limite sinistro sono diversi.
\end{itemize}
Poiché $\lim_{x\to 1^-} M(x) \neq \lim_{x\to 1^+} M(x)$, la funzione $M$ in $x_0 = 1$ ha una \emph{discontinuità di prima specie}, con salto
\[
\lim_{x\to 1^+} M(x) - \lim_{x\to 1^-} M(x) = 0 - 1 = -1 .
\]
\end{esempio}

\begin{esempio}[Discontinuità di seconda specie]
Sia
\[
h(x) = \begin{cases} \dfrac{1}{x} & x \neq 0 \\ 0 & x = 0 .\end{cases}
\]
Poiché $\displaystyle\lim_{x\to 0^+}\frac{1}{x} = +\infty$ e $\displaystyle\lim_{x\to 0^-}\frac{1}{x} = -\infty$, la funzione in $x_0 = 0$ ha una discontinuità di seconda specie.

Sia
\[
\ell(x) = \begin{cases} \dfrac{1}{x} & x > 0 \\ 0 & x \le 0 .\end{cases}
\]
Si ha $\displaystyle\lim_{x\to 0^+}\ell(x) = \lim_{x\to 0^+}\frac{1}{x} = +\infty$ e $\displaystyle\lim_{x\to 0^-}\ell(x) = \lim_{x\to 0^-} 0 = 0$: la funzione $\ell$ in $x_0 = 0$ ha una discontinuità di seconda specie (basta che uno dei due limiti sia infinito).
\end{esempio}

\end{document}
